\documentclass[a4paper, serif, 10pt]{moderncv}

%% moderncv
% color and style
\moderncvcolor{black}
\moderncvstyle{banking}

%% page layout/margins
\usepackage[top=2.5cm, bottom=2.5cm, left=1.5cm, right=1.5cm]{geometry}

\nopagenumbers{}

%% Babel config for Indonesian (Bahasa Indonesia) language
% ref:
% - https://latex3.github.io/babel/guides/locale-indonesian.html
\usepackage[indonesian]{babel}

%% fontspec
\usepackage{fontspec}

\IfFontExistsTF{Linux Libertine}{
  \setmainfont[Ligatures={TeX, Common}, Numbers=OldStyle]{Linux Libertine}
}{}
\IfFontExistsTF{Linux Libertine O}{
  \setmainfont[Ligatures={TeX, Common}, Numbers=OldStyle]{Linux Libertine O}
}{}

%% CTeX
% CJK support, used to show CJK characters in the resume
%
% - heading=false: disable ctex heading
% - punct=kaiming: use kaiming punctuations styles for CJK
% - scheme=plain: use plain scheme, do not override \normalsize font size
% - space=auto: space settings for CJK characters
%
% CJK fonts are resolved per locale language, e.g. Noto Serif/Sans CJK JP
% for Japanese resumes, with fallback to the next region if not installed
%
% ref:
% - http://ctan.mirrorcatalogs.com/language/chinese/ctex/ctex.pdf
\usepackage[UTF8, heading=false, punct=kaiming, scheme=plain, space=auto]{ctex}

\IfFontExistsTF{Noto Serif CJK SC}{
  \setCJKmainfont{Noto Serif CJK SC}
}{}
\IfFontExistsTF{Noto Sans CJK SC}{
  \setCJKsansfont{Noto Sans CJK SC}
}{}

%% line spacing
\usepackage{setspace}
\setstretch{1.125}

%% URL styling - use normal text instead of monospace
\urlstyle{same}

%% Patch \httplink and \httpslink to handle full URLs
% The original moderncv macros always prepend http:// or https:// to URLs.
% This causes issues when the URL already has a protocol (e.g., https://example.com)
% resulting in malformed URLs like https://https://example.com.
% This patch checks if the URL already contains "://" and uses it directly if so.
% Note: We use \str_set:Nx (x-expansion) to expand macros like \@homepage first.
\ExplSyntaxOn
\str_new:N \g_yamlresume_url_str

\RenewDocumentCommand{\httpslink}{O{}m}{
  \str_set:Nx \g_yamlresume_url_str {#2}
  \str_if_in:NnTF \g_yamlresume_url_str {://}
    { \url{#2} }
    { \url{https://#2} }
}

\RenewDocumentCommand{\httplink}{O{}m}{
  \str_set:Nx \g_yamlresume_url_str {#2}
  \str_if_in:NnTF \g_yamlresume_url_str {://}
    { \url{#2} }
    { \url{http://#2} }
}
\ExplSyntaxOff

\name{Rizky Pratama Nugraha}{}
\title{Analis Keuangan Senior | Pemodelan Finansial \& Valuasi}
\phone[mobile]{+62 812 3456 7890}
\email{rizky.pratama@keuangannusantara.co.id}
\homepage{https://www.rizkypratama.id}

\address{Jl. Jenderal Sudirman Kav. 52-53, Jakarta Selatan, DKI Jakarta, Indonesia, 12190}{}{}

\extrainfo{{\small \faLinkedin}\ \href{https://www.linkedin.com/in/rizky-pratama-keuangan}{@rizky-pratama-keuangan} {} {} {} • {} {} {} 
{\small \faMedium}\ \href{https://medium.com/@rizkypratama.finance}{@rizkypratama.finance}}

\begin{document}

\maketitle

\section{Informasi Dasar}

\cvline{}{Analis keuangan dengan pengalaman lebih dari 8 tahun dalam analisis investasi, pemodelan finansial, dan manajemen risiko di sektor perbankan dan pasar modal.
%
\begin{itemize}
\item Terampil membangun model valuasi DCF, komparasi perusahaan, dan analisis sensitivitas untuk mendukung keputusan investasi.
\item Berpengalaman menyusun laporan keuangan manajemen, proyeksi kas, dan rekomendasi portofolio bagi manajemen senior.
\item Kuat dalam analisis data menggunakan Excel lanjutan, SQL, Python, dan Tableau; memahami regulasi OJK dan standar PSAK.
\item Terbiasa bekerja lintas fungsi, memimpin riset industri, serta menyampaikan insight kepada pemangku kepentingan.
\end{itemize}}
\section{Pendidikan}

\cventry{Sep 2014–Jul 2016}
        {Magister, Manajemen Keuangan, Nilai: 3.85}
        {Universitas Cakrawala Bisnis}
        {\url{https://www.ucb.ac.id}}
        {}
        {\begin{itemize}
\item Memperdalam teknik valuasi perusahaan, analisis risiko, dan strategi portofolio.
\item Menyelesaikan tesis tentang dampak suku bunga terhadap kinerja sektor perbankan.
\end{itemize}
\textbf{Mata Kuliah}: Analisis Investasi, Manajemen Risiko Keuangan, Pasar Modal dan Portofolio, Akuntansi Manajemen Lanjutan, Ekonomi Moneter}

\cventry{Sep 2010–Jul 2014}
        {Sarjana, Akuntansi, Nilai: 3.72}
        {Universitas Nusantara Ekonomi}
        {\url{https://www.une.ac.id}}
        {}
        {\begin{itemize}
\item Membangun dasar kuat dalam laporan keuangan, penganggaran, dan analisis rasio.
\item Aktif dalam organisasi kemahasiswaan bidang pasar modal.
\end{itemize}
\textbf{Mata Kuliah}: Akuntansi Keuangan, Manajemen Keuangan, Statistik Bisnis, Ekonomi Mikro dan Makro, Audit Internal}
\section{Pengalaman Kerja}

\cventry{Mar 2021–Sekarang}
        {Analis Keuangan Senior}
        {PT Sentosa Kapital Investama}
        {\url{https://www.sentosa-kapital.co.id}}
        {}
        {\begin{itemize}
\item Memimpin analisis investasi untuk portofolio ekuitas dan obligasi senilai Rp2,4 triliun.
\item Membangun model DCF, LBO, dan komparasi perusahaan untuk mendukung uji tuntas.
\item Menyusun laporan kinerja bulanan dan rekomendasi alokasi aset kepada direksi.
\item Mengawasi dua analis junior dan mengembangkan template pemodelan standar.
\end{itemize}
\textbf{Kata Kunci}: Analisis Investasi, Pemodelan DCF, Uji Tuntas, Valuasi Perusahaan, Riset Industri}

\cventry{Jun 2017–Feb 2021}
        {Analis Keuangan}
        {Bank Nusantara Mandiri}
        {\url{https://www.banknusantaramandiri.co.id}}
        {}
        {\begin{itemize}
\item Menganalisis kinerja kredit korporasi dan menyusun proyeksi arus kas debitur.
\item Menilai kelayakan pembiayaan dengan analisis rasio, sensitivitas, dan uji stres.
\item Menyiapkan paket laporan manajemen untuk komite kredit dan regulator.
\item Berkolaborasi dengan tim risiko dan kepatuhan untuk memastikan kesesuaian regulasi OJK.
\end{itemize}
\textbf{Kata Kunci}: Analisis Kredit, Proyeksi Kas, Manajemen Risiko, Regulasi OJK, Pelaporan Manajemen}

\cventry{Agu 2016–Mei 2017}
        {Analis Keuangan Junior}
        {PT Analitika Pasar Modal}
        {\url{https://www.analitikapasarmodal.co.id}}
        {}
        {\begin{itemize}
\item Mengumpulkan dan membersihkan data pasar untuk riset sektor konsumen dan perbankan.
\item Menyusun ringkasan harian pasar modal dan catatan analis untuk tim investasi.
\item Membantu pembuatan model keuangan dasar serta presentasi internal.
\end{itemize}
\textbf{Kata Kunci}: Riset Pasar, Data Pasar, Model Keuangan, Pasar Modal, Analisis Sektor}
\section{Bahasa}

\cvline{Indonesia}{Bahasa Ibu / Bilingual}
\cvline{Inggris}{Tingkat Profesional Penuh \hfill \textbf{Kata Kunci}: IELTS 7.5, Presentasi Bisnis}
\section{Keahlian}

\cvline{Analisis Keuangan}{Ahli \hfill \textbf{Kata Kunci}: Laporan Keuangan, Analisis Rasio, Proyeksi Kas, Anggaran, Valuasi}
\cvline{Pemodelan Finansial}{Mahir \hfill \textbf{Kata Kunci}: DCF, LBO, Excel Lanjutan, Analisis Sensitivitas, Komparasi Perusahaan}
\cvline{Analitik Data}{Mahir \hfill \textbf{Kata Kunci}: SQL, Python, Tableau, Power BI, Bloomberg Terminal}
\cvline{Manajemen Risiko}{Menengah \hfill \textbf{Kata Kunci}: Uji Stres, Analisis Kredit, Regulasi OJK, PSAK, Kepatuhan}
\cvline{Komunikasi Bisnis}{Mahir \hfill \textbf{Kata Kunci}: Presentasi Manajemen, Penulisan Laporan, Kolaborasi Lintas Fungsi}
\section{Penghargaan}

\cventry{Des 2023}
        {PT Sentosa Kapital Investama}
        {Analis Keuangan Terbaik}
        {}
        {}
        {Diberikan atas kontribusi dalam penyusunan model valuasi portofolio dan rekomendasi investasi yang meningkatkan imbal hasil portofolio inti sebesar 12\% secara tahunan.}
\section{Sertifikat}

\cventry{Nov 2022}
        {Corporate Finance Institute}
        {Financial Modeling \& Valuation Analyst (FMVA)}
        {\url{https://corporatefinanceinstitute.com/certifications/fmva/}}
        {}
        {}

\cventry{Agu 2020}
        {CFA Institute}
        {Certified Financial Analyst (CFA) Level I}
        {\url{https://www.cfainstitute.org/programs/cfa-program}}
        {}
        {}
\section{Publikasi}

\cventry{Jun 2022}
        {Dampak Suku Bunga terhadap Kinerja Sektor Perbankan Indonesia}
        {Jurnal Riset Keuangan Nusantara}
        {\url{https://www.jrkn.ac.id/artikel/dampak-suku-bunga-perbankan}}
        {}
        {\begin{itemize}
\item Menganalisis hubungan antara perubahan suku bunga acuan dan profitabilitas bank umum.
\item Menemukan bahwa kualitas kredit dan likuiditas memoderasi dampak suku bunga terhadap margin bunga bersih.
\item Memberikan rekomendasi kebijakan bagi manajemen risiko perbankan.
\end{itemize}}
\section{Proyek}

\cventry{Jan 2023–Agu 2023}
        {Dasbor internal untuk memantau KPI keuangan, arus kas, dan eksposur risiko portofolio.}
        {Dasbor Pemantauan Kinerja Keuangan}
        {\url{https://www.sentosa-kapital.co.id/proyek/dashboard-keuangan}}
        {}
        {\begin{itemize}
\item Merancang metrik utama dan tata kelola data untuk dasbor manajemen.
\item Mengotomatiskan laporan bulanan menggunakan SQL, Python, dan Power BI.
\item Memangkas waktu penyusunan laporan dari 5 hari menjadi 1 hari.
\end{itemize}
\textbf{Kata Kunci}: Dasbor, Power BI, Otomasi Laporan, KPI Keuangan, Tata Kelola Data}

\cventry{Mar 2022–Des 2022}
        {Model valuasi untuk menilai peluang investasi di sektor energi terbarukan.}
        {Model Valuasi Sektor Energi Terbarukan}
        {\url{https://www.sentosa-kapital.co.id/proyek/valuasi-energi}}
        {}
        {\begin{itemize}
\item Membangun model DCF dan skenario sensitivitas harga listrik serta biaya modal.
\item Menyusun komparasi perusahaan dan analisis risiko regulasi.
\item Menyajikan rekomendasi investasi kepada komite investasi.
\end{itemize}
\textbf{Kata Kunci}: DCF, Energi Terbarukan, Analisis Skenario, Risiko Regulasi, Komite Investasi}
\section{Minat}

\cvline{Pasar Modal}{Analisis Saham, Obligasi, Reksa Dana, Makroekonomi}
\cvline{Pengembangan Diri}{Berbicara Publik, Menulis, Visualisasi Data}
\section{Kegiatan Sukarela}

\cventry{Jan 2018–Des 2024}
        {Relawan Edukasi Keuangan}
        {Yayasan Literasi Keuangan Indonesia}
        {\url{https://www.ylki.or.id}}
        {}
        {\begin{itemize}
\item Memberikan pelatihan literasi keuangan kepada pelajar dan pelaku UMKM.
\item Menyusun modul sederhana tentang anggaran, tabungan, dan investasi aman.
\item Membantu kampanye pencegahan investasi ilegal bersama komunitas lokal.
\end{itemize}}
    
\end{document}